% Methods: the vertical datum chain (docs/DATUMS.md). Design: D-069, D-070, D-073, D-074;
% PREREG section 5 and Amendment 11. Every number comes from \metric{...}.

\subsection*{Vertical datums}
\emph{Working datum.} All heights were compared on the reference of the DAHITI series: normal heights, ellipsoidal altimeter heights minus a quasi-geoid model\citep{schwatke2015}, EIGEN-6C4 in the current data version\citep{dahiti_wl,forste2014}. Level bins, waterlines, topography and hindcast all use it. Levels are reported in the Baltic Height System; a constant offset between the two systems applies equally to levels and to heights brought to the DAHITI reference, so it cancels in every comparison here and is estimated only for reporting (pre-registration section~5).

\emph{Relief grids.} For the total error of the hindcast audit, the grids were read as heights above the EGM2008 geoid\citep{pavlis2012}. ETOPO~2022 states this reference in its file metadata\citep{macferrin2025}; GEBCO and ETOPO1 describe their heights as referred to sea level\citep{gebco2026,amante2009}, and the SRTM15+ distribution files state none\citep{tozer2019}, so for these grids the EGM2008 reading is an assumption. Each grid was brought to the DAHITI reference by adding, cell by cell, $N_{\mathrm{EGM2008}} - N_{\mathrm{EIGEN\text{-}6C4}}$, computed from the two tide-free spherical-harmonic models with the calculation service of the International Centre for Global Earth Models\citep{ince2019}. Over the domain the difference ranges from \metric{datum.dN_min_m} to \metric{datum.dN_max_m}~m (standard deviation \metric{datum.dN_sd_m}~m). The widely distributed realisation of EGM2008, which ETOPO~2022's geoid layer reproduces, carries a further constant of \metric{datum.nga_offset_m}~m\todo{cite: NGA EGM2008 distribution documentation}; the total error uses it, and the total error without it is a sensitivity. The shape error fits its own constant and depends on neither choice.

\emph{Lidar.} ICESat-2 ATL06 and ATL08 heights (Release~007) are tide-free ellipsoidal heights on WGS84 in the ITRF2020 frame\citep{atl06_v7,atl08_v7,is2_compguide_r007}. ATL06 carries the tide-free EGM2008 geoid height $N$, and heights above the geoid are $H = h - N$; for ATL08, which has no geoid, $N$ was interpolated in time along the same beam of the same orbit's ATL06 granule from samples at most \metric{params.geoid_max_gap_s}~s apart, and unmatched points were dropped and counted. A constant $c$, the median of the topography minus $H$ over training-block points, with a block-bootstrap uncertainty, brings $H$ to the DAHITI reference\todo{check before evaluation: caspian.validation applies only the constant $c$; the per-cell $N_{\mathrm{EGM2008}} - N_{\mathrm{EIGEN\text{-}6C4}}$ difference used for the relief grids is not applied to the lidar heights, so its spatial part remains in the residuals}. DGFI does not state the tide system of DAHITI heights; if it is mean-tide, a small north--south tilt remains and is stated as an error. GEDI heights (pre-registration Amendment~2) are treated as ICESat-2. GLO-30 heights are referenced to EGM2008\citep{copdem_handbook_i5}, and its two geocells in the domain without TanDEM-X measurements were excluded.

\emph{SWOT.} SWOT elevations are heights above a geoid that the files do not name, after removal of solid-earth, load and pole tides. Its difference from the DAHITI reference is constant in time within a block and is removed with each block's mean departure in the setup model.
